\section{Acceleration: Why a Small Lag Grows}
\label{sec:accel}

If technologies now spread faster, why worry about lags? Because what
matters for income is not whether a technology has \emph{arrived} but how
intensively it is \emph{used}. \textcite{comin2010} showed that adoption lags
have shortened dramatically across successive technologies;
\textcite{comin2018} found that while adoption lags converged, the intensity of use in poor
countries relative to rich ones diverged---and that the productivity trends
implied by these two patterns resemble the actual divergence of world incomes
over the past two centuries.

\begin{figure}[t]
\centering
\begin{tikzpicture}
\begin{axis}[edge, height=0.58\linewidth,
    xmin=0, xmax=72, xlabel={Median adoption lag behind the first five adopters, years},
    ymode=log, ymin=4, ymax=60, ytick={5,10,20,50}, yticklabels={5$\times$,10$\times$,20$\times$,50$\times$},
    ylabel={Intensity gap, richest quarter vs poorest half},
    legend pos=north east]
  \addplot[only marks, mark=*, mark size=2.4pt, mark options={fill=haze,draw=slate}, only if={era}{industrial}]
    table[x=lag,y=gap_ratio] {\datadir/adoption.csv};
  \addlegendentry{Industrial era (frontier before 1915)}
  \addplot[only marks, mark=square*, mark size=2.2pt, mark options={fill=teal,draw=white}, only if={era}{mid}]
    table[x=lag,y=gap_ratio] {\datadir/adoption.csv};
  \addlegendentry{Mid-century (1915--1960)}
  \addplot[only marks, mark=diamond*, mark size=3pt, mark options={fill=ember,draw=white}, only if={era}{digital}]
    table[x=lag,y=gap_ratio] {\datadir/adoption.csv};
  \addlegendentry{Digital (after 1975)}
  \addplot[labelled=label] table[x=lag,y=gap_ratio] {\datadir/adoption.csv};
\end{axis}
\end{tikzpicture}
\caption{\takeaway{Technologies arrive faster but are used less
  equally.} For \nTechs{} technologies, the typical country adopted
  industrial-era inventions decades after the frontier, but digital ones within
  about a decade. Yet in 2000 the richest quarter of countries used computers and
  mobile phones \gapComputer{} and \gapCell{} times as intensively per person
  as the poorest half, against \gapOldMedian{} times for the median
  pre-1950 technology. Digital technologies were young in 2000, which
  inflates their gaps; see text.
  \source{CHAT \parencite{comin2009chat}; author's calculations following
  \textcite{comin2010,comin2018}.}}
\label{fig:adoption}
\end{figure}


\Cref{fig:adoption} reproduces that pattern descriptively with the CHAT
dataset \parencite{comin2009chat}. For each of \nTechs{} technologies, the
horizontal axis shows how long the typical country took to adopt it after the
first five adopters; the vertical axis shows how many times more intensively
the richest quarter of countries used it, per person, than the poorest half.
Lags shrank from \lagRail{} years for passenger railways and
\lagTelephone{} for the telephone to \lagCell{} for mobile phones and
\lagComputer{} for computers. But intensity gaps grew: \gapRail-fold for
railways and \gapTV-fold for television, \gapCell-fold for mobile phones and
\gapComputer-fold for computers in 2000. Across technologies, shorter lags
go with \emph{larger} gaps (correlation \corrLagGap).

An honest caveat: in 2000 mobile phones and computers were young, and gaps
during early diffusion are naturally large; some of the difference will close,
as it has for internet access (\cref{fig:internet}). The structural point
survives the caveat. When the frontier moves fast, a country that adopts
late and uses lightly is chasing a moving target, and each new layer of
technology builds on the last. Being a few years and a few multiples behind
in broadband becomes being further behind in cloud services, in data-driven
logistics, in AI.

\paragraph{Thin layers.} Three further mechanisms turn small lags into
lasting gaps in small, remote places.
\emph{Language}: digital services, voice assistants and now AI models are
built first, and best, for languages with many speakers and much text; a
language with a hundred thousand speakers is served last, if at all.
\emph{Food systems}: once local agriculture and supply chains have been
displaced by imports, the knowledge and infrastructure to rebuild them
decays. \emph{Thin institutions}: a ministry with a handful of staff must
cover the same range of international standards, negotiations and
regulations as one with thousands. Each new technology---data protection,
digital trade rules, AI governance---adds a portfolio but not a person.
